% R=5 first-law attempt: occupancy-stable hop probe dies.
\documentclass[11pt]{article}
\usepackage[margin=1.05in]{geometry}
\usepackage{amsmath,amssymb,amsthm}
\usepackage[colorlinks=true,linkcolor=blue,citecolor=blue,urlcolor=blue]{hyperref}

\newtheorem*{admission}{Admission}

\title{\textbf{At Larger Radius the Occupancy-Stable\\
First-Law Probe Dies}}
\author{Robert W.\ McGwier\\
\small CTO, Cohere Technology Group\\
\small GitHub: \texttt{n4hy}}
\date{15 August 2026 --- Version 1}

\begin{document}
\maketitle

\begin{abstract}
\noindent
Paper~29 left C4 (CHM parabola versus linear $R-r$) as a tie on
a ball of radius $4$. A larger ball is the next place those
envelopes can split. It is not a step toward the Planck length.

$N=14$, $R=5$, thirty locked surface hops. At $\varepsilon=0.01$,
$27$ of $30$ flip $n_{\mathrm{occ}}$. At $\varepsilon=0.002$,
$25$ of $30$ still flip. The five survivors are not a census.
C4 is not scored.

The last valid uniqueness test remains Paper~29: $294$ hops,
zero flips, a $0.002$ residual tie. This note records an
instrument halt, not a nearer ultraviolet.
\end{abstract}

\section{Why it dies}

The unperturbed sea has $n_{\mathrm{occ}}=1370$ on $14^3=2744$
sites. Half filling is $1372$. A surface hop of size $0.002$
moves a near-zero mode across the Fermi level. That is not
linear Clausius. It is a level crossing.

\begin{admission}
I went to larger $R$ to split the parabola from a linear weight.
The probe that was stable at $R=4$ is not stable at $R=5$. I did
not call that the Planck scale, and I did not score five leftover
hops as a theorem.
\end{admission}

\end{document}
